\documentclass[12pt,a4paper]{exam}
\usepackage[utf8]{inputenc}
\usepackage{amsmath,amssymb,amsthm}
\usepackage[margin=1in]{geometry}
\usepackage{graphicx}
\usepackage[colorlinks=true]{hyperref}
\pagestyle{headandfoot}
\firstpageheader{MIT OpenCourseWare}{}{\textbf{EXTREMELY HARD -- MIT LEVEL}}
\runningheader{MIT OpenCourseWare}{}{\textbf{EXTREMELY HARD -- MIT LEVEL}}
\firstpagefooter{}{Page \thepage}{}
\runningfooter{}{Page \thepage}{}

\begin{document}

\begin{center}
{\LARGE \textbf{MIT OpenCourseWare}}\\14.06\\Advanced Macroeconomics
\vspace{0.5cm}
{14.06} --- {Advanced Macroeconomics}
\vspace{0.3cm}
May 2026
\vspace{0.3cm}
\textbf{EXTREMELY HARD -- MIT LEVEL}
\end{center}

\begin{center}
\textbf{Time Limit: 3 hours}
\end{center}

\begin{center}
\textbf{Total Marks: 100}
\end{center}

\vspace{0.5cm}

\begin{questions}

\question[2]
Construct an original proof-level problem involving DSGE Models for 14.06 (Advanced Macroeconomics) and provide a complete solution. Your problem should be at the level of a graduate qualifying exam.

\vspace{0.3cm}

\question[2]
Construct an original proof-level problem involving Real Business Cycles for 14.06 (Advanced Macroeconomics) and provide a complete solution. Your problem should be at the level of a graduate qualifying exam.

\vspace{0.3cm}

\question[5]
Prove the fundamental theorem concerning New Keynesian Economics in the context of 14.06 (Advanced Macroeconomics). State the theorem precisely, provide a complete proof, and discuss three important corollaries.

\vspace{0.3cm}

\question[3]
Analyze the limitations of standard approaches to Optimal Monetary Policy when applied to edge cases in 14.06. Construct explicit counterexamples showing where intuitive reasoning fails.

\vspace{0.3cm}

\question[2]
Prove the fundamental theorem concerning Fiscal Theory in the context of 14.06 (Advanced Macroeconomics). State the theorem precisely, provide a complete proof, and discuss three important corollaries.

\vspace{0.3cm}

\question[6]
Analyze the limitations of standard approaches to DSGE Models when applied to edge cases in 14.06. Construct explicit counterexamples showing where intuitive reasoning fails.

\vspace{0.3cm}

\question[5]
Analyze the limitations of standard approaches to Real Business Cycles when applied to edge cases in 14.06. Construct explicit counterexamples showing where intuitive reasoning fails.

\vspace{0.3cm}

\question[4]
Prove the fundamental theorem concerning New Keynesian Economics in the context of 14.06 (Advanced Macroeconomics). State the theorem precisely, provide a complete proof, and discuss three important corollaries.

\vspace{0.3cm}

\question[10]
Construct an original proof-level problem involving Optimal Monetary Policy for 14.06 (Advanced Macroeconomics) and provide a complete solution. Your problem should be at the level of a graduate qualifying exam.

\vspace{0.3cm}

\question[3]
Construct an original proof-level problem involving Fiscal Theory for 14.06 (Advanced Macroeconomics) and provide a complete solution. Your problem should be at the level of a graduate qualifying exam.

\vspace{0.3cm}

\question[3]
Construct an original proof-level problem involving DSGE Models for 14.06 (Advanced Macroeconomics) and provide a complete solution. Your problem should be at the level of a graduate qualifying exam.

\vspace{0.3cm}

\question[4]
Prove the fundamental theorem concerning Real Business Cycles in the context of 14.06 (Advanced Macroeconomics). State the theorem precisely, provide a complete proof, and discuss three important corollaries.

\vspace{0.3cm}

\question[3]
Construct an original proof-level problem involving New Keynesian Economics for 14.06 (Advanced Macroeconomics) and provide a complete solution. Your problem should be at the level of a graduate qualifying exam.

\vspace{0.3cm}

\question[5]
Analyze the limitations of standard approaches to Optimal Monetary Policy when applied to edge cases in 14.06. Construct explicit counterexamples showing where intuitive reasoning fails.

\vspace{0.3cm}

\question[2]
Analyze the limitations of standard approaches to Fiscal Theory when applied to edge cases in 14.06. Construct explicit counterexamples showing where intuitive reasoning fails.

\vspace{0.3cm}

\question[3]
Construct an original proof-level problem involving DSGE Models for 14.06 (Advanced Macroeconomics) and provide a complete solution. Your problem should be at the level of a graduate qualifying exam.

\vspace{0.3cm}

\question[5]
Prove the fundamental theorem concerning Real Business Cycles in the context of 14.06 (Advanced Macroeconomics). State the theorem precisely, provide a complete proof, and discuss three important corollaries.

\vspace{0.3cm}

\question[4]
Prove the fundamental theorem concerning New Keynesian Economics in the context of 14.06 (Advanced Macroeconomics). State the theorem precisely, provide a complete proof, and discuss three important corollaries.

\vspace{0.3cm}

\question[3]
Construct an original proof-level problem involving Optimal Monetary Policy for 14.06 (Advanced Macroeconomics) and provide a complete solution. Your problem should be at the level of a graduate qualifying exam.

\vspace{0.3cm}

\question[4]
Prove the fundamental theorem concerning Fiscal Theory in the context of 14.06 (Advanced Macroeconomics). State the theorem precisely, provide a complete proof, and discuss three important corollaries.

\vspace{0.3cm}

\question[2]
Analyze the limitations of standard approaches to DSGE Models when applied to edge cases in 14.06. Construct explicit counterexamples showing where intuitive reasoning fails.

\vspace{0.3cm}

\question[7]
Analyze the limitations of standard approaches to Real Business Cycles when applied to edge cases in 14.06. Construct explicit counterexamples showing where intuitive reasoning fails.

\vspace{0.3cm}

\question[3]
Construct an original proof-level problem involving New Keynesian Economics for 14.06 (Advanced Macroeconomics) and provide a complete solution. Your problem should be at the level of a graduate qualifying exam.

\vspace{0.3cm}

\question[4]
Construct an original proof-level problem involving Optimal Monetary Policy for 14.06 (Advanced Macroeconomics) and provide a complete solution. Your problem should be at the level of a graduate qualifying exam.

\vspace{0.3cm}

\question[6]
Construct an original proof-level problem involving Fiscal Theory for 14.06 (Advanced Macroeconomics) and provide a complete solution. Your problem should be at the level of a graduate qualifying exam.

\vspace{0.3cm}

\end{questions}

\newpage
\section*{Solutions}

\textbf{Solution 1:}

Solution approach: First recall the definitions and key theorems related to DSGE Models from 14.06. Then apply the appropriate techniques to construct a rigorous argument. The solution must be self-contained and reference only the material from the course.

\vspace{0.5cm}

\textbf{Solution 2:}

Solution approach: First recall the definitions and key theorems related to Real Business Cycles from 14.06. Then apply the appropriate techniques to construct a rigorous argument. The solution must be self-contained and reference only the material from the course.

\vspace{0.5cm}

\textbf{Solution 3:}

The proof follows from the definitions and properties established in 14.06. The key insight is to apply the relevant theorem and verify the hypotheses hold. Detailed solution depends on the specific formulation of the theorem.

\vspace{0.5cm}

\textbf{Solution 4:}

The edge case analysis reveals the subtle assumptions underlying standard treatments of Optimal Monetary Policy. By constructing explicit counterexamples, we see that the usual hypotheses cannot be weakened without loss of the theorem's conclusion.

\vspace{0.5cm}

\textbf{Solution 5:}

The proof follows from the definitions and properties established in 14.06. The key insight is to apply the relevant theorem and verify the hypotheses hold. Detailed solution depends on the specific formulation of the theorem.

\vspace{0.5cm}

\textbf{Solution 6:}

The edge case analysis reveals the subtle assumptions underlying standard treatments of DSGE Models. By constructing explicit counterexamples, we see that the usual hypotheses cannot be weakened without loss of the theorem's conclusion.

\vspace{0.5cm}

\textbf{Solution 7:}

The edge case analysis reveals the subtle assumptions underlying standard treatments of Real Business Cycles. By constructing explicit counterexamples, we see that the usual hypotheses cannot be weakened without loss of the theorem's conclusion.

\vspace{0.5cm}

\textbf{Solution 8:}

The proof follows from the definitions and properties established in 14.06. The key insight is to apply the relevant theorem and verify the hypotheses hold. Detailed solution depends on the specific formulation of the theorem.

\vspace{0.5cm}

\textbf{Solution 9:}

Solution approach: First recall the definitions and key theorems related to Optimal Monetary Policy from 14.06. Then apply the appropriate techniques to construct a rigorous argument. The solution must be self-contained and reference only the material from the course.

\vspace{0.5cm}

\textbf{Solution 10:}

Solution approach: First recall the definitions and key theorems related to Fiscal Theory from 14.06. Then apply the appropriate techniques to construct a rigorous argument. The solution must be self-contained and reference only the material from the course.

\vspace{0.5cm}

\textbf{Solution 11:}

Solution approach: First recall the definitions and key theorems related to DSGE Models from 14.06. Then apply the appropriate techniques to construct a rigorous argument. The solution must be self-contained and reference only the material from the course.

\vspace{0.5cm}

\textbf{Solution 12:}

The proof follows from the definitions and properties established in 14.06. The key insight is to apply the relevant theorem and verify the hypotheses hold. Detailed solution depends on the specific formulation of the theorem.

\vspace{0.5cm}

\textbf{Solution 13:}

Solution approach: First recall the definitions and key theorems related to New Keynesian Economics from 14.06. Then apply the appropriate techniques to construct a rigorous argument. The solution must be self-contained and reference only the material from the course.

\vspace{0.5cm}

\textbf{Solution 14:}

The edge case analysis reveals the subtle assumptions underlying standard treatments of Optimal Monetary Policy. By constructing explicit counterexamples, we see that the usual hypotheses cannot be weakened without loss of the theorem's conclusion.

\vspace{0.5cm}

\textbf{Solution 15:}

The edge case analysis reveals the subtle assumptions underlying standard treatments of Fiscal Theory. By constructing explicit counterexamples, we see that the usual hypotheses cannot be weakened without loss of the theorem's conclusion.

\vspace{0.5cm}

\textbf{Solution 16:}

Solution approach: First recall the definitions and key theorems related to DSGE Models from 14.06. Then apply the appropriate techniques to construct a rigorous argument. The solution must be self-contained and reference only the material from the course.

\vspace{0.5cm}

\textbf{Solution 17:}

The proof follows from the definitions and properties established in 14.06. The key insight is to apply the relevant theorem and verify the hypotheses hold. Detailed solution depends on the specific formulation of the theorem.

\vspace{0.5cm}

\textbf{Solution 18:}

The proof follows from the definitions and properties established in 14.06. The key insight is to apply the relevant theorem and verify the hypotheses hold. Detailed solution depends on the specific formulation of the theorem.

\vspace{0.5cm}

\textbf{Solution 19:}

Solution approach: First recall the definitions and key theorems related to Optimal Monetary Policy from 14.06. Then apply the appropriate techniques to construct a rigorous argument. The solution must be self-contained and reference only the material from the course.

\vspace{0.5cm}

\textbf{Solution 20:}

The proof follows from the definitions and properties established in 14.06. The key insight is to apply the relevant theorem and verify the hypotheses hold. Detailed solution depends on the specific formulation of the theorem.

\vspace{0.5cm}

\textbf{Solution 21:}

The edge case analysis reveals the subtle assumptions underlying standard treatments of DSGE Models. By constructing explicit counterexamples, we see that the usual hypotheses cannot be weakened without loss of the theorem's conclusion.

\vspace{0.5cm}

\textbf{Solution 22:}

The edge case analysis reveals the subtle assumptions underlying standard treatments of Real Business Cycles. By constructing explicit counterexamples, we see that the usual hypotheses cannot be weakened without loss of the theorem's conclusion.

\vspace{0.5cm}

\textbf{Solution 23:}

Solution approach: First recall the definitions and key theorems related to New Keynesian Economics from 14.06. Then apply the appropriate techniques to construct a rigorous argument. The solution must be self-contained and reference only the material from the course.

\vspace{0.5cm}

\textbf{Solution 24:}

Solution approach: First recall the definitions and key theorems related to Optimal Monetary Policy from 14.06. Then apply the appropriate techniques to construct a rigorous argument. The solution must be self-contained and reference only the material from the course.

\vspace{0.5cm}

\textbf{Solution 25:}

Solution approach: First recall the definitions and key theorems related to Fiscal Theory from 14.06. Then apply the appropriate techniques to construct a rigorous argument. The solution must be self-contained and reference only the material from the course.

\vspace{0.5cm}

\end{document}